\BookRecto
\begingroup
\renewcommand{\runningtitle}{Pronunciation}
\hyphenpenalty=10000
\exhyphenpenalty=10000
\emergencystretch=3em
\newcommand{\reserve}[1]{\par\penalty-100\begingroup\dimen0=\pagegoal\advance\dimen0 by-\pagetotal\ifdim\dimen0<#1\newpage\fi\endgroup}
\newcommand{\itempara}[2]{\par\addvspace{3pt}\begingroup\raggedright\noindent\hangindent=1.8em\hangafter=1\makebox[1.8em][l]{#1.}#2\par\endgroup}
% book-source-body:begin

\addfontfeatures{Ligatures=TeXOff}
\raggedright
\booktitle{\BookComponent{pronunciation-learner}{A guide to Old English pronunciation}GUIDE}{A guide to Old English pronunciation}

% mdparagraph:P001:begin
This guide gives a practical pronunciation for reading the course. English examples refer to a particular sound in an ordinary American pronunciation, not to the whole word. Where English has no equivalent, the instructions describe the position of the tongue and lips. Every Old English word is glossed; its meaning need not be memorized.
% mdparagraph:P001:end

\reserve{4\baselineskip}
% mdparagraph:P002:begin
\section*{1. Letters and length}\BookHeading{nav-pronunciation-learner-0001}{subsection}{1. Letters and length}{1}\BookLabel{bk-pronunciation-learner-156b3d5ac4b3}
% mdparagraph:P002:end

% mdparagraph:P003:begin
\textbf{Æ æ}, called \emph{ash}, is one vowel letter. \textbf{Þ þ}, \emph{thorn}, and \textbf{Ð ð}, \emph{eth}, represent the sounds written \textbf{th} in English \emph{thin} and \emph{this}. Either letter can represent either sound. This course normally uses \textbf{ð}. Cook's hooked \textbf{ę} and \textbf{ǫ} have the sounds of \textbf{e} and \textbf{o}; the hooks record their history, not additional sounds.[1]
% mdparagraph:P003:end

% mdparagraph:P004:begin
A \textbf{macron}, the line above a vowel, marks length: \textbf{ā, ǣ, ē, ī, ō, ū, ȳ}. Hold a long vowel appreciably longer than a short one. Long does not mean the English letter name, and a macron does not mark stress. Thus \textbf{ac}, “but”, has a short vowel; \textbf{āc}, “oak”, has a long one. Both end with the sound of \textbf{k}.
% mdparagraph:P004:end

% mdparagraph:P005:begin
In a long two-letter vowel spelling, the macron stands on the first letter: \textbf{ēa, ēo, īe}. The two letters belong to one syllable. Preserve the printed marks; a following pair of consonants does not turn an unmarked vowel into a long one.
% mdparagraph:P005:end

\reserve{4\baselineskip}
% mdparagraph:P006:begin
\section*{2. Simple vowels}\BookHeading{nav-pronunciation-learner-0002}{subsection}{2. Simple vowels}{1}\BookLabel{bk-pronunciation-learner-c95193d444e8}
% mdparagraph:P006:end

% mdparagraph:P007:begin
Use these short and long values. Some of the practical English approximations differ in quality as well as duration.[2]
% mdparagraph:P007:end

% mdtable:T01:begin
% course-table-style: 2026-10
% course-panel:T01-P01:begin
\par\begingroup
\fontsize{10.5}{13.2}\selectfont
\setlength{\LTleft}{0pt}\setlength{\LTright}{\fill}
\noindent
\begin{tabular}{@{}>{\raggedright\arraybackslash}p{15.099mm}>{\raggedright\arraybackslash}p{52.318mm}>{\raggedright\arraybackslash}p{50.959mm}@{}}
\toprule
\textbf{Letters} & \textbf{Short vowel} & \textbf{Long vowel} \\
\midrule
a, ā & The vowel of \emph{father}, brief: \textbf{ac}, “but”. & Hold the same vowel: \textbf{āc}, “oak”. \\
æ, ǣ & The vowel of \emph{cat}: \textbf{sæt}, “sat”. & Hold that same vowel: \textbf{sǣd}, “seed”. \\
e, ē & The vowel of \emph{bet}: \textbf{męte}, “food”. & Hold the beginning of \emph{day} without its final glide: \textbf{hē}, “he”. \\
i, ī & The vowel of \emph{bit}: \textbf{is}, “is”. & The vowel of \emph{see}, held: \textbf{īs}, “ice”. \\
o, ō & A brief rounded vowel, described below: \textbf{god}, “god”. & Hold the beginning of \emph{go} without its final glide: \textbf{gōd}, “good”. \\
u, ū & The vowel of \emph{foot}: \textbf{full}, “full”. & The vowel of \emph{food}, held: \textbf{fūl}, “vile”. \\
y, ȳ & The tongue position of \emph{bit} with rounded lips: \textbf{wyrt}, “herb; plant”. & The tongue position of \emph{see} with rounded lips, held: \textbf{fȳr}, “fire”. \\
\bottomrule
\end{tabular}
\par\endgroup
% course-panel:T01-P01:end
% mdtable:T01:end

% mdparagraph:P008:begin
For \textbf{o}, the vowel of \emph{caught} is a useful starting point if you distinguish \emph{cot} from \emph{caught}. If they sound alike, begin with the first part of \emph{go}, lower the jaw a little while keeping the lips rounded, and make the vowel brief. Keep \textbf{o} distinct from the unrounded \emph{father} vowel. For \textbf{ō}, keep the jaw and tongue steady: do not finish with an \textbf{oo} movement.
% mdparagraph:P008:end

% mdparagraph:P009:begin
For \textbf{ē}, likewise stop before the glide at the end of \emph{day}. Do not pronounce \textbf{ǣ} like American \emph{care}: it is the \emph{cat} vowel held longer, without an \textbf{r} sound.
% mdparagraph:P009:end

% mdparagraph:P010:begin
In a \textbf{stressed} syllable, short \textbf{a} before \textbf{m} or \textbf{n} takes the more rounded short-\textbf{o} quality just described. Use this in \textbf{mann}, “man; human being”, and \textbf{lang}, “long”. It is not the vowel of ordinary American \emph{man}. This rule does not require the same rounding in every unstressed \textbf{-an} ending.
% mdparagraph:P010:end

% mdparagraph:P011:begin
For \textbf{y} and \textbf{ȳ}, keep the tongue forward while rounding the lips as for \textbf{oo} in \emph{food}. For short \textbf{y}, begin with the tongue position of \emph{bit}; for long \textbf{ȳ}, use \emph{see}. Rounding the lips must not pull the tongue back into an ordinary \textbf{oo} vowel. These sounds have no usual American English equivalent. Here \textbf{y} is a vowel, not the consonant at the beginning of \emph{yes}.
% mdparagraph:P011:end

\reserve{4\baselineskip}
% mdparagraph:P012:begin
\section*{3. Diphthongs}\BookHeading{nav-pronunciation-learner-0003}{subsection}{3. Diphthongs}{1}\BookLabel{bk-pronunciation-learner-1212ca9f676b}
% mdparagraph:P012:end

% mdparagraph:P013:begin
A \emph{diphthong} is a vowel that moves from one position toward another within one syllable. For these spellings, make one smooth movement toward \textbf{oo}. Do not give the second letter a separate beat.[3]
% mdparagraph:P013:end

% mdtable:T02:begin
% course-table-style: 2026-10
% course-panel:T02-P01:begin
\par\begingroup
\fontsize{10.5}{13.2}\selectfont
\setlength{\LTleft}{0pt}\setlength{\LTright}{\fill}
\noindent
\begin{tabular}{@{}>{\raggedright\arraybackslash}p{19.374mm}>{\raggedright\arraybackslash}p{63.45mm}>{\raggedright\arraybackslash}p{35.552mm}@{}}
\toprule
\textbf{Spellings} & \textbf{Movement} & \textbf{Supplied examples} \\
\midrule
ea, ēa & Start with the vowel of \emph{cat}, then round toward \textbf{oo}. & \textbf{earm}, “arm”; \textbf{ēare}, “ear”. \\
eo, ēo & Start with the steady first vowel of \emph{day}, then round toward \textbf{oo}. & \textbf{eorl}, “man”; \textbf{fēoll}, “fell”. \\
ie, īe & Start with the vowel of \emph{see}, then round toward \textbf{oo}. & \textbf{giefan}, “give”; \textbf{hīeran}, “hear”. \\
\bottomrule
\end{tabular}
\par\endgroup
% course-panel:T02-P01:end
% mdtable:T02:end

% mdparagraph:P014:begin
In \textbf{ēa, ēo, īe}, hold the first part longer before making the movement. The \textbf{eo} of \textbf{eorl} is short; the \textbf{ēo} of \textbf{fēoll} is long. The word \textbf{ēare} has two syllables: its long \textbf{ēa}, followed by a light final \textbf{e}. The opening \textbf{g} of \textbf{giefan} has the consonant sound of \textbf{y} in \emph{yes}.
% mdparagraph:P014:end

% mdparagraph:P015:begin
There are word-specific exceptions. In \textbf{geong}, “young”, \textbf{e} does not supply a vowel sound: the initial \textbf{g} is like \textbf{y} in \emph{yes}, followed by short rounded \textbf{o}. Its \textbf{ng} has the sound heard in \emph{finger}, including the \textbf{g}. Do not give this word the ordinary \textbf{eo} glide or copy the vowel of modern \emph{young}.
% mdparagraph:P015:end

\reserve{4\baselineskip}
% mdparagraph:P016:begin
\section*{4. Consonants}\BookHeading{nav-pronunciation-learner-0004}{subsection}{4. Consonants}{1}
% mdparagraph:P016:end

% mdparagraph:P017:begin
Read \textbf{b, d, l, m, n, p, t, w} as at the beginning of \emph{boy, dog, lip, man, no, pen, top, win}. Use a pronounced American \textbf{r}, as in \emph{run}, as a practical approximation. The remaining letters and combinations need the following distinctions.[4]
% mdparagraph:P017:end

\reserve{4\baselineskip}
% mdparagraph:P018:begin
\subsection*{c, g, and combinations}\BookHeading{nav-pronunciation-learner-0005}{subsubsection}{c, g, and combinations}{0}\BookLabel{bk-pronunciation-learner-16b43725bb81}
% mdparagraph:P018:end

% mdparagraph:P019:begin
Work in the following order. A supplied pronunciation for a particular word or form takes precedence over a general pattern.
% mdparagraph:P019:end

% mdparagraph:P020:begin
\textbf{1. Read a combination first.} Do not apply a rule for single \textbf{c} or \textbf{g} to one letter of these combinations.
% mdparagraph:P020:end

% mdtable:T03:begin
% course-table-style: 2026-10
% course-panel:T03-P01:begin
\par\begingroup
\fontsize{10.5}{13.2}\selectfont
\setlength{\LTleft}{0pt}\setlength{\LTright}{\fill}
\noindent
\begin{tabular}{@{}>{\raggedright\arraybackslash}p{15.233mm}>{\raggedright\arraybackslash}p{35.876mm}>{\raggedright\arraybackslash}p{67.267mm}@{}}
\toprule
\textbf{Spelling} & \textbf{Ordinary sound} & \textbf{Model and supplied exception} \\
\midrule
sc & \textbf{sh} as in \emph{ship}, including before a back vowel. & \textbf{scip}, “ship”, has one \textbf{sh} sound. The exception \textbf{frosc}, “frog”, ends with \textbf{s} followed by \textbf{k}. \\
cg; usually gg & The final sound of \emph{bridge}. & \textbf{ecg}, “edge”. The exception \textbf{docga}, “dog”, has a prolonged hard \textbf{g}. \\
ng & As in \emph{finger}, with audible \textbf{g}. & \textbf{lang}, “long”. In the supplied exception \textbf{sengan}, “singe”, \textbf{ng} has the sound at the end of \emph{singe}. \\
\bottomrule
\end{tabular}
\par\endgroup
% course-panel:T03-P01:end
% mdtable:T03:end

% mdparagraph:P021:begin
For the rules below, the \textbf{back-vowel spellings} are \textbf{a/\allowbreak{}ā, o/\allowbreak{}ō, u/\allowbreak{}ū}. The relevant \textbf{front-vowel spellings} are \textbf{e/\allowbreak{}ē, i/\allowbreak{}ī, æ/\allowbreak{}ǣ}; the diphthongs \textbf{ea/\allowbreak{}ēa, eo/\allowbreak{}ēo, ie/\allowbreak{}īe} also begin at the front of the mouth. Although \textbf{y/\allowbreak{}ȳ} are front vowels in articulation, their spelling is not a reliable instruction to give \textbf{c} or \textbf{g} a front sound. Use the supplied pronunciation in such words.
% mdparagraph:P021:end

% mdparagraph:P022:begin
\textbf{2. At the beginning of a word, decide c from what immediately follows.}\par\nobreak
% mdparagraph:P022:end

% mdtable:T04:begin
% course-table-style: 2026-10
% course-panel:T04-P01:begin
\par\begingroup
\fontsize{10.5}{13.2}\selectfont
\setlength{\LTleft}{0pt}\setlength{\LTright}{\fill}
\noindent
\begin{tabular}{@{}>{\raggedright\arraybackslash}p{49.566mm}>{\raggedright\arraybackslash}p{27.057mm}>{\raggedright\arraybackslash}p{41.752mm}@{}}
\toprule
\textbf{Following sound or spelling} & \textbf{Ordinary c sound} & \textbf{Model} \\
\midrule
A consonant, or a/\allowbreak{}ā, o/\allowbreak{}ō, u/\allowbreak{}ū & \textbf{k} as in \emph{kit} & \textbf{cnīf}, “knife”; \textbf{cuman}, “come”. \\
e/\allowbreak{}ē, i/\allowbreak{}ī, æ/\allowbreak{}ǣ, or a front diphthong & \textbf{ch} as in \emph{chin} & \textbf{cealf}, “calf”, begins with \textbf{ch}. \\
\bottomrule
\end{tabular}
\par\endgroup
% course-panel:T04-P01:end
% mdtable:T04:end

% mdparagraph:P023:begin
These are useful defaults, with lexical exceptions. \textbf{Cēlan}, “cool”, begins with \textbf{k} despite \textbf{ē}. Both \textbf{c} letters in \textbf{cyrice}, “church”, are \textbf{ch}: learn its initial sound with the word, rather than deducing it from \textbf{y}. Keep all three syllables, including the final \textbf{e}.
% mdparagraph:P023:end

% mdparagraph:P024:begin
\textbf{3. Inside or at the end of a word, use the neighbours of c.} Before another consonant, the ordinary sound is \textbf{k}. Immediately after \textbf{i/\allowbreak{}ī}, use \textbf{ch} when \textbf{c} ends the word or is followed by a vowel other than \textbf{a/\allowbreak{}ā, o/\allowbreak{}ō, u/\allowbreak{}ū}. The second \textbf{c} of \textbf{cyrice} follows this pattern. A following back vowel instead gives the ordinary \textbf{k} pattern.
% mdparagraph:P024:end

% mdparagraph:P025:begin
Compare \textbf{dīc}, “ditch; dike”, with \textbf{dīcas}, “ditches; dikes”. Final \textbf{c} after \textbf{ī} has \textbf{ch}; before the \textbf{a} of the longer form it has \textbf{k}. An ending can therefore change the consonant sound. In other positions \textbf{k} is the ordinary starting point, but some words need a supplied pronunciation. Do not extend the after-\textbf{i/\allowbreak{}ī} rule to every front vowel.
% mdparagraph:P025:end

% mdparagraph:P026:begin
Learn \textbf{sēcan}, “seek”, as a supplied exception: its \textbf{c} is \textbf{ch} as in \emph{chin}, despite the following \textbf{a}. The back-vowel pattern does not override a word's stated pronunciation.
% mdparagraph:P026:end

% mdparagraph:P027:begin
\textbf{4. For single g, first decide whether it has the consonant sound y in \emph{yes}.}\par\nobreak
% mdparagraph:P027:end

% mdtable:T05:begin
% course-table-style: 2026-10
% course-panel:T05-P01:begin
\par\begingroup
\fontsize{10.5}{13.2}\selectfont
\setlength{\LTleft}{0pt}\setlength{\LTright}{\fill}
\noindent
\begin{tabular}{@{}>{\raggedright\arraybackslash}p{46.0mm}>{\raggedright\arraybackslash}p{20.0mm}>{\raggedright\arraybackslash}p{52.376mm}@{}}
\toprule
\textbf{Position} & \textbf{Ordinary g sound} & \textbf{Model} \\
\midrule
At the beginning, immediately before a front vowel or front diphthong & \textbf{y} as in \emph{yes} & \textbf{giefan}, “give”. The prefix \textbf{ge-} also begins with this sound. \\
After a front vowel, at the end or when no back vowel immediately follows & \textbf{y} as in \emph{yes} & \textbf{dæg}, “day”, ends with \textbf{y}. A following consonant does not by itself require hard \textbf{g}. \\
\bottomrule
\end{tabular}
\par\endgroup
% course-panel:T05-P01:end
% mdtable:T05:end

% mdparagraph:P028:begin
Otherwise use the non-\textbf{y} branch below, unless a word has a supplied exceptional pronunciation. At the beginning, a consonant between \textbf{g} and a later front vowel prevents the first row from applying: look at the immediate neighbour. The \textbf{g} rule after a front vowel is broader than the \textbf{c} rule after \textbf{i/\allowbreak{}ī}; keep them distinct. The earlier caution about \textbf{y/\allowbreak{}ȳ} still applies.
% mdparagraph:P028:end

% mdparagraph:P029:begin
\textbf{5. For non-y g, decide between stopped and continuous g.}\par\nobreak
% mdparagraph:P029:end

% mdtable:T06:begin
% course-table-style: 2026-10
% course-panel:T06-P01:begin
\par\begingroup
\fontsize{10.5}{13.2}\selectfont
\setlength{\LTleft}{0pt}\setlength{\LTright}{\fill}
\noindent
\begin{tabular}{@{}>{\raggedright\arraybackslash}p{47.364mm}>{\raggedright\arraybackslash}p{37.919mm}>{\raggedright\arraybackslash}p{33.092mm}@{}}
\toprule
\textbf{Position} & \textbf{Sound} & \textbf{Model} \\
\midrule
At the beginning; after n; in a supplied hard doubled consonant & Stopped \textbf{g}, as in \emph{go} & \textbf{god}, “god”; \textbf{lang}, “long”; \textbf{docga}, “dog”. \\
Elsewhere & The continuous voiced back sound described below & \textbf{fugol}, “bird”: single \textbf{g} between back vowels. \\
\bottomrule
\end{tabular}
\par\endgroup
% course-panel:T06-P01:end
% mdtable:T06:end

% mdparagraph:P030:begin
For\BookPoint{bk-pronunciation-learner-b426b0db6cb3} the continuous \textbf{g}, put the back of the tongue near the position for hard \textbf{g}, but leave a narrow gap. Let air pass continuously while the throat vibrates: the sound can be held, unlike the stopped \textbf{g} of \emph{go}. Its phonetic sign is \textbf{[ɣ]}. Do not omit it or substitute \textbf{y}. First selecting the \textbf{y} or non-\textbf{y} branch prevents this rule from changing the final sound of \textbf{dæg} into continuous \textbf{g}.
% mdparagraph:P030:end

% mdparagraph:P031:begin
These patterns let you decide regular new examples. Where a word or form has an exceptional pronunciation, use its stated sound; spelling alone does not settle every case. The practice below supplies any required exception. No historical derivation is required.
% mdparagraph:P031:end

\reserve{4\baselineskip}
% mdparagraph:P032:begin
\subsection*{h and sounded clusters}\BookHeading{nav-pronunciation-learner-0006}{subsubsection}{h and sounded clusters}{0}\BookLabel{bk-pronunciation-learner-fce112ac8644}
% mdparagraph:P032:end

% mdparagraph:P033:begin
At the beginning of a syllable, \textbf{h} is like English \textbf{h} in \emph{hat}. Later in a syllable it has a stronger, continuous sound made by narrowing the passage above the tongue.[5]
% mdparagraph:P033:end

% mdparagraph:P034:begin
In \textbf{brōhte}, “brought”, raise the back of the tongue toward the position for \textbf{k}, leaving a gap for breath. Keep the throat unvoiced. This back sound is \textbf{[x]}. Both it and the following \textbf{t} are heard.
% mdparagraph:P034:end

% mdparagraph:P035:begin
In \textbf{niht}, “night”, the preceding \textbf{i} brings the tongue farther forward. Raise the middle of the tongue toward the roof of the mouth and let breath pass without voice. This front sound is \textbf{[ç]}. Keep the \textbf{t} too: the Old English word does not have modern English's silent \textbf{gh}. Front vowels generally favour this front articulation; back vowels favour the back one.
% mdparagraph:P035:end

% mdparagraph:P036:begin
Sound the initials in \textbf{hl, hn, hr, hw, wr, wl, cn}. In \textbf{hlāf}, “loaf”, begin with audible breath before \textbf{l}; in \textbf{hwæl}, “whale”, keep \textbf{h} with \textbf{w}. In \textbf{wrītan}, “write”, pronounce \textbf{w} and \textbf{r}; in \textbf{cnīf}, “knife”, pronounce \textbf{k} and \textbf{n}. Join the consonants without putting a vowel between them. The spelling \textbf{cw} has \textbf{k} followed by \textbf{w}. Read \textbf{x} and \textbf{hs} as \textbf{k} followed by \textbf{s}; \textbf{weax}, “wax”, ends with that pair.
% mdparagraph:P036:end

\reserve{4\baselineskip}
% mdparagraph:P037:begin
\section*{5. f, s, and th: voice}\BookHeading{nav-pronunciation-learner-0007}{subsection}{5. f, s, and th: voice}{1}\BookLabel{bk-pronunciation-learner-74b76bb98627}
% mdparagraph:P037:end

% mdparagraph:P038:begin
Compare \emph{fan} with \emph{van}, \emph{sip} with \emph{zip}, and \emph{thin} with \emph{this}. The initial consonant of the second word has voice: the throat vibrates. Vowels and consonants such as \textbf{l, m, n, r} also carry voice.[6]
% mdparagraph:P038:end

% mdparagraph:P039:begin
The ordinary rule for single \textbf{f, s, þ/\allowbreak{}ð} is:
% mdparagraph:P039:end

\begin{itemize}
% mdparagraph:P040:begin
\item At the beginning or end of a word, or beside a voiceless consonant, use the sounds of \emph{fan, sip, thin}.
% mdparagraph:P040:end

% mdparagraph:P041:begin
\item Inside a word, between voiced sounds, use the sounds of \emph{van, zip, this}.
% mdparagraph:P041:end

\end{itemize}
% mdtable:T07:begin
% course-table-style: 2026-10
% course-panel:T07-P01:begin
\par\begingroup
\fontsize{10.5}{13.2}\selectfont
\setlength{\LTleft}{0pt}\setlength{\LTright}{\fill}
\noindent
\begin{tabular}{@{}>{\raggedright\arraybackslash}p{26.306mm}>{\raggedright\arraybackslash}p{26.306mm}>{\raggedright\arraybackslash}p{65.764mm}@{}}
\toprule
\textbf{Word} & \textbf{Meaning} & \textbf{Sound to use} \\
\midrule
cnīf & knife & Final \textbf{f} as in \emph{fan}. \\
cnīfas & knives & \textbf{f} as in \emph{van}, between vowels; final \textbf{s} as in \emph{sip}. \\
rīsan & rise & \textbf{s} as in \emph{zip}, between vowels. \\
hūs & house & Final \textbf{s} as in \emph{sip}. \\
ōðer & other & \textbf{ð} as in \emph{this}, between vowels. \\
ðā & then & Initial \textbf{ð} as in \emph{thin}. \\
\bottomrule
\end{tabular}
\par\endgroup
% course-panel:T07-P01:end
% mdtable:T07:end

% mdparagraph:P042:begin
Thus \textbf{is}, “is”, ends with the \textbf{s} of \emph{sip}, not the \textbf{z} of its modern English equivalent. The letter \textbf{ð} does not by itself tell you to use the voiced sound.
% mdparagraph:P042:end

% mdparagraph:P043:begin
Stress and the boundaries between word parts can preserve exceptions. In \textbf{hǣlðu}, “health”, and \textbf{forðām}, “therefore”, use the unvoiced \textbf{th} of \emph{thin}. Learn these as supplied word-specific readings; the surrounding voiced sounds do not override them.
% mdparagraph:P043:end

\reserve{4\baselineskip}
% mdparagraph:P044:begin
\section*{6. Duration, syllables, and stress}\BookHeading{nav-pronunciation-learner-0008}{subsection}{6. Duration, syllables, and stress}{1}\BookLabel{bk-pronunciation-learner-a18005f2eb82}
% mdparagraph:P044:end

% mdparagraph:P045:begin
Keep vowel length separate from consonant length. In \textbf{full}, “full”, \textbf{u} is short; \textbf{fūl}, “vile”, has long \textbf{ū}. A doubled consonant at the end of a word normally has ordinary single-consonant duration: the final \textbf{ll} of \textbf{full} does not require extra holding.
% mdparagraph:P045:end

% mdparagraph:P046:begin
Between vowels, a doubled consonant has extra duration, like the \textbf{n} across the join in \emph{pen-knife}. Do not insert a vowel between its two letters.[7]
% mdparagraph:P046:end

% mdparagraph:P047:begin
With a doubled stop, hold the closure longer. In \textbf{habban}, “have”, keep the lips closed for \textbf{bb}, then release into the following vowel. There are two syllables. Do not add another vowel or another syllable to pronounce the doubling.
% mdparagraph:P047:end

% mdparagraph:P048:begin
Written unstressed vowels remain audible and distinguishable. In \textbf{męte}, “food”, there are two syllables, and final \textbf{e} is a light \emph{bet} vowel. In \textbf{fugol}, “bird”, there are also two syllables; its \textbf{o} is heard. Do not reduce every ending to the \textbf{uh} at the beginning of modern \emph{about}.
% mdparagraph:P048:end

% mdparagraph:P049:begin
The main stress normally falls on the first syllable, as in \textbf{habban}, “have”. Prefixes complicate this: \textbf{ge-} is unstressed, \textbf{be-} and \textbf{for-} usually are, and a prefixed verb generally stresses its root. When \textbf{ge-} precedes a one-syllable root, the word has two syllables and the root receives the stress. A supplied division identifies the prefix where needed. A macron indicates duration regardless of which syllable is stressed.
% mdparagraph:P049:end

% mdparagraph:P050:begin
For these careful reading exercises, observe printed length. In connected speech, unstressed function words and prefixes could be shorter than their conventional written macrons suggest. This does not authorize adding, removing, or moving marks in the course text.
% mdparagraph:P050:end

\reserve{4\baselineskip}
% mdparagraph:P051:begin
\section*{7. Practice}\BookHeading{nav-pronunciation-learner-0009}{subsection}{7. Practice}{1}
% mdparagraph:P051:end

% mdparagraph:P052:begin
The following tasks are new teaching material. Meanings are supplied throughout. Apply the guide's sound patterns to the new words; no translation or grammatical analysis is required. Give an English sound cue or the taught mouth instruction, not a phonetic transcription.[9]
% mdparagraph:P052:end

\reserve{4\baselineskip}
% mdparagraph:P053:begin
\subsection*{A. Vowels and length}\BookHeading{nav-pronunciation-learner-0010}{subsubsection}{A. Vowels and length}{0}
% mdparagraph:P053:end

% mdparagraph:P054:begin
Read each word aloud. Keep duration separate from stress, and give each diphthong one syllable.
% mdparagraph:P054:end

% item: PRON-P01a
% mdparagraph:P055:begin
\textbf{A1.} \textbf{for}, “for”; \textbf{fōr}, “went”. What changes in the vowel? Give its sound and duration in each word.
% mdparagraph:P055:end

% item: PRON-P01b
% mdparagraph:P056:begin
\textbf{A2.} \textbf{węnde}, “went”; \textbf{wēnde}, “thought; supposed”. Explain the vowel contrast. Does the hook under \textbf{ę} require a sound different from ordinary short \textbf{e}? Keep the final \textbf{e} audible in both words.
% mdparagraph:P056:end

% item: PRON-P01c
% mdparagraph:P057:begin
\textbf{A3.} \textbf{dǣd}, “deed”; \textbf{wīd}, “wide”; \textbf{hȳdan}, “hide; conceal”. Give the sound of each marked vowel. Explain how the lips distinguish \textbf{ȳ} from \textbf{ī}, then read all three words.
% mdparagraph:P057:end

% item: PRON-P01d
% mdparagraph:P058:begin
\textbf{A4.} \textbf{eall}, “all”; \textbf{scēap}, “sheep”; \textbf{heorte}, “heart”; \textbf{frēond}, “friend”; \textbf{scield}, “shield”; \textbf{gīet}, “yet; still”. For each diphthong, give its starting sound and movement, and say whether the first part is short or long. Count the syllables in each word without giving the diphthong two beats.
% mdparagraph:P058:end

\reserve{4\baselineskip}
% mdparagraph:P059:begin
\subsection*{B. Consonants and unfamiliar sounds}\BookHeading{nav-pronunciation-learner-0011}{subsubsection}{B. Consonants and unfamiliar sounds}{0}
% mdparagraph:P059:end

% mdparagraph:P060:begin
Choose the sound of the indicated letters and explain which position or neighbouring sound selects it. Then read the whole word. Check combinations before single letters.
% mdparagraph:P060:end

% item: PRON-P02a
% mdparagraph:P061:begin
\textbf{B1.} \textbf{c} in \textbf{cēosað}, “choose” (a plural form). Which initial-\textbf{c} pattern applies?
% mdparagraph:P061:end

% item: PRON-P02b
% mdparagraph:P062:begin
\textbf{B2.} \textbf{c} in \textbf{micel}, “large; great”, and \textbf{micles}, “large; great” (a declined form). Do the two forms have the same consonant sound? Explain from the immediate neighbours of \textbf{c}.
% mdparagraph:P062:end

% item: PRON-P02c
% mdparagraph:P063:begin
\textbf{B3.} \textbf{g} in \textbf{gān}, “go; walk”, and \textbf{glæd}, “glad”. Does the \textbf{æ} of the second word make its initial \textbf{g} a \textbf{y} sound? Apply both stages of the \textbf{g} decision.
% mdparagraph:P063:end

% item: PRON-P02d
% mdparagraph:P064:begin
\textbf{B4.} \textbf{g} in \textbf{fæger}, “beautiful; fair”, and \textbf{regn}, “rain”. Decide from the sounds before and after \textbf{g}. Does the following \textbf{n} require stopped \textbf{g} in the second word?
% mdparagraph:P064:end

% item: PRON-P02e
% mdparagraph:P065:begin
\textbf{B5.} \textbf{g} in \textbf{āgan}, “have; possess”. Choose among \textbf{y}, stopped \textbf{g}, and continuous \textbf{g}, and explain both decisions. Show with the tongue and airflow how to make the selected sound.
% mdparagraph:P065:end

% item: PRON-P02f
% mdparagraph:P066:begin
\textbf{B6.} \textbf{cg} in \textbf{brycg}, “bridge”, and \textbf{ng} in \textbf{singan}, “sing”. Which English cue fits each combination? Keep the short rounded vowel of \textbf{brycg} distinct from ordinary English \textbf{i}.
% mdparagraph:P066:end

% item: PRON-P02g
% mdparagraph:P067:begin
\textbf{B7.} \textbf{h} in \textbf{riht}, “right; direct”, and \textbf{dohtor}, “daughter”. What changes in the tongue position, and which vowel causes the change? Keep the following \textbf{t} audible in both words.
% mdparagraph:P067:end

% item: PRON-P02h
% mdparagraph:P068:begin
\textbf{B8.} \textbf{sc} in \textbf{scūfan}, “thrust”, and \textbf{cw} in \textbf{cwæð}, “said”. Decide whether each spelling gives one consonant sound or a sequence. Does the back vowel \textbf{ū} require a separate \textbf{k} in the first word?
% mdparagraph:P068:end

\reserve{4\baselineskip}
% mdparagraph:P069:begin
\subsection*{C. Voice}\BookHeading{nav-pronunciation-learner-0012}{subsubsection}{C. Voice}{0}
% mdparagraph:P069:end

% mdparagraph:P070:begin
State whether the indicated sound is voiced or unvoiced. Give its English sound cue and explain the context or supplied exception.
% mdparagraph:P070:end

% item: PRON-P03a
% mdparagraph:P071:begin
\textbf{C1.} \textbf{f} in \textbf{ǣfre}, “ever; always”. Identify the sounds on both sides; they need not both be vowels.
% mdparagraph:P071:end

% item: PRON-P03b
% mdparagraph:P072:begin
\textbf{C2.} Both \textbf{f} letters in \textbf{fīf}, “five”. Decide each from its position rather than copying the pronunciation of modern English \emph{five}.
% mdparagraph:P072:end

% item: PRON-P03c
% mdparagraph:P073:begin
\textbf{C3.} \textbf{s} and final \textbf{ð} in \textbf{cēosað}, “choose” (a plural form). Does one voicing decision apply to both letters?
% mdparagraph:P073:end

% item: PRON-P03d
% mdparagraph:P074:begin
\textbf{C4.} \textbf{f} and \textbf{ð} in \textbf{hæfð}, “has”. Begin with the word-final consonant, then use its voicing to decide the preceding \textbf{f}.
% mdparagraph:P074:end

% item: PRON-P03e
% mdparagraph:P075:begin
\textbf{C5.} \textbf{ð} in \textbf{cwæð}, “said”. Would replacing its spelling with thorn, \textbf{þ}, change the sound? Explain the role of position and the choice of letter.
% mdparagraph:P075:end

% item: PRON-P03f
% mdparagraph:P076:begin
\textbf{C6.} Compare \textbf{ð} in \textbf{hǣlðu}, “health”, with \textbf{f} in \textbf{ǣfre}, “ever; always”. The unvoiced \textbf{th} in \textbf{hǣlðu} is supplied lexical information from §5. Explain why that fact takes precedence over the ordinary between-voiced-sounds rule, while the ordinary rule applies to \textbf{ǣfre}. This is an exception comparison, not a request to deduce the exception from its spelling.
% mdparagraph:P076:end

\reserve{4\baselineskip}
% mdparagraph:P077:begin
\subsection*{D. Syllables and stress}\BookHeading{nav-pronunciation-learner-0013}{subsubsection}{D. Syllables and stress}{0}
% mdparagraph:P077:end

% mdparagraph:P078:begin
Give the number of syllables and identify the stressed one. Identify any doubled consonant and explain how it is pronounced.
% mdparagraph:P078:end

% item: PRON-P04a
% mdparagraph:P079:begin
\textbf{D1.} \textbf{sittan}, “sit”. Does prolonging \textbf{tt} lengthen \textbf{i}, or add a syllable?
% mdparagraph:P079:end

% item: PRON-P04b
% mdparagraph:P080:begin
\textbf{D2.} \textbf{sunne}, “sun”. Keep the last vowel distinct. Explain the duration of \textbf{nn} and whether it changes the syllable count.
% mdparagraph:P080:end

% item: PRON-P04c
% mdparagraph:P081:begin
\textbf{D3.} \textbf{gewrit}, “letter; written document”; division \textbf{ge-writ}. Apply the prefix rule to stress, and pronounce both \textbf{w} and \textbf{r} without inserting another vowel.
% mdparagraph:P081:end

% item: PRON-P04d
% mdparagraph:P082:begin
\textbf{D4.} \textbf{forlǣtan}, “abandon; forsake”; division \textbf{for-lǣtan}. This is a prefixed verb. Locate its stress from the prefix rule, and explain what the macron tells you separately.
% mdparagraph:P082:end

\reserve{4\baselineskip}
% mdparagraph:P083:begin
\subsection*{E. A sentence adapted from Cook}\BookHeading{nav-pronunciation-learner-0014}{subsubsection}{E. A sentence adapted from Cook}{0}
% mdparagraph:P083:end

% mdparagraph:P084:begin
This sentence adapts the opening of Cook's Creation reading.[8]
% mdparagraph:P084:end

% mdparagraph:P085:begin
\textbf{God gescēop heofonan and eorðan.}\par\nobreak
% mdparagraph:P085:end

% mdparagraph:P086:begin
“God created heaven and earth.”
% mdparagraph:P086:end

% mdtable:T08:begin
% course-table-style: 2026-10
% course-panel:T08-P01:begin
\par\begingroup
\fontsize{10.5}{13.2}\selectfont
\setlength{\LTleft}{0pt}\setlength{\LTright}{\fill}
\noindent
\begin{tabular}{@{}>{\raggedright\arraybackslash}p{56.106mm}>{\raggedright\arraybackslash}p{65.082mm}@{}}
\toprule
\textbf{Word} & \textbf{Supplied meaning} \\
\midrule
God & God \\
gescēop & created \\
heofonan & heaven \\
and & and \\
eorðan & earth \\
\bottomrule
\end{tabular}
\par\endgroup
% course-panel:T08-P01:end
% mdtable:T08:end

% mdparagraph:P087:begin
The verb has the division \textbf{ge-scēop}, with prefix \textbf{ge-}. The sentence's meaning and grammatical forms are supplied; apply the pronunciation rules to read it aloud, then answer these sound questions.
% mdparagraph:P087:end

% item: PRON-P05a
% mdparagraph:P088:begin
\textbf{E1.} Which sounds do \textbf{g} and \textbf{sc} have in \textbf{gescēop}, “created”? Explain the rules you use. Count its syllables and place the stress, keeping \textbf{ēo} one glide.
% mdparagraph:P088:end

% item: PRON-P05b
% mdparagraph:P089:begin
\textbf{E2.} What sound does \textbf{f} have in \textbf{heofonan}, “heaven”? Identify its neighbours and explain your decision. Count the word's syllables.
% mdparagraph:P089:end

% item: PRON-P05c
% mdparagraph:P090:begin
\textbf{E3.} What sound does \textbf{ð} have in \textbf{eorðan}, “earth”? Identify its neighbours and explain why the ordinary rule applies even though one neighbour is a consonant.
% mdparagraph:P090:end

\reserve{4\baselineskip}
% mdparagraph:P091:begin
\section*{Source notes}\BookHeading{nav-pronunciation-learner-0015}{subsection}{Source notes}{1}
% mdparagraph:P091:end

% mdparagraph:P092:begin
\itempara{1}{Cook, \emph{A First Book in Old English}, second edition (1895), §§2–4, pp.5–6 /\allowbreak{} PDF21–22; Fulk, \emph{An Introductory Grammar of Old English}, corrected open-access reprint (2020), §§15–16, pp.10–12 /\allowbreak{} PDF22–24. Fulk's teaching dots are omitted in the selected examples; \textbf{hǣlðu, forðām} replace his thorn with the course's eth. These declared spelling changes do not change their selected sounds.}
% mdparagraph:P092:end

% mdparagraph:P093:begin
\itempara{2}{Fulk §§12,16,19, pp.10–11,13–14 /\allowbreak{} PDF22–23,25–26, permits the traditional short-vowel approximations used here. Held \emph{cat} for \textbf{ǣ} follows Fulk. Cook §5, pp.6–7 /\allowbreak{} PDF22–23, also offers practical approximations, but some of his English cues suit other accents. American cue adjustments are teaching aids.}
% mdparagraph:P093:end

% mdparagraph:P094:begin
\itempara{3}{Fulk §§20–21, p.14 /\allowbreak{} PDF26; \textbf{giefan}, §135a, p.103 /\allowbreak{} PDF115. The rounded diphthongs follow Fulk's declared convention. Cook's §5 table, p.7 /\allowbreak{} PDF23, gives different component vowels; the two descriptions are not interchangeable.}
% mdparagraph:P094:end

% mdparagraph:P095:begin
\itempara{4}{Fulk §§16,18,135a–d, pp.11–13,103–104 /\allowbreak{} PDF23–25,115–116; Cook §§7,10, pp.8–9 /\allowbreak{} PDF24–25. \textbf{Cealf, dīc, dīcas} follow Fulk §135b, with his pronunciation overpoints omitted; Cook's vocabulary, p.283 /\allowbreak{} PDF299, also gives \textbf{dīc}. For \textbf{sēcan}, Cook's vocabulary, p.311 /\allowbreak{} PDF327, supplies the form and meaning as a headword variant; Fulk §135b supports the \textbf{ch} value with its explicitly prefixed infinitive example. The conditions for noninitial \textbf{c} are narrower than those for \textbf{g}. Earlier vowel changes, later endings and analogy explain some exceptions to the practical spelling patterns; the sound-history appendix treats these matters separately.}
% mdparagraph:P095:end

% mdparagraph:P096:begin
\itempara{5}{Cook §§7,10, pp.8–9 /\allowbreak{} PDF24–25; Fulk §§9,16,18, pp.8,11–13 /\allowbreak{} PDF20,23–25. American \textbf{r} is a practical approximation. The \textbf{k+s} reading of \textbf{x/\allowbreak{}hs} follows Cook's practical value; Fulk also records an earlier h-like sound before \textbf{s}. Neither source requires a trill in these practical descriptions.}
% mdparagraph:P096:end

% mdparagraph:P097:begin
\itempara{6}{Fulk §17, p.12 /\allowbreak{} PDF24, supplies the voicing rule and exceptions used here. Cook §8, pp.8–9 /\allowbreak{} PDF24–25, gives a broader syllable-final voicing rule. This guide selects Fulk's account.}
% mdparagraph:P097:end

% mdparagraph:P098:begin
\itempara{7}{Cook §11, pp.9–10 /\allowbreak{} PDF25–26, distinguishes prolonged medial doubling from final double spellings, normally read with single-consonant duration. See also Cook §23, pp.18–19 /\allowbreak{} PDF34–35, and Fulk §§17,19,22–23, pp.12–15 /\allowbreak{} PDF24–27. Printed vowel length and stress are distinct; connected-speech shortening does not require changing the edited spelling.}
% mdparagraph:P098:end

% mdparagraph:P099:begin
\itempara{8}{Adapted from Cook, Reader I, p.123 /\allowbreak{} PDF139, first prose line: remove \textbf{On anginne}, move \textbf{God} before the verb, and omit the three apparatus cues. The retained words are not respelled. The original sentence and notes remain in the source transcription. The full meaning and grammatical forms are supplied here; the sound questions apply the guide's rules. The separate \textbf{ge-scēop} division identifies a prefix and does not respell the reading.}
% mdparagraph:P099:end

% mdparagraph:P100:begin
\itempara{9}{The new practice forms are checked against Fulk §§16–24,135, pp.11–15,103–104 /\allowbreak{} PDF23–27,115–116; Cook §§4–5,53,121, pp.6–7,37,74 /\allowbreak{} PDF22–23,53,90, Reader I p.123 /\allowbreak{} PDF139, and vocabulary pp.283,288–289,296,309–312 /\allowbreak{} PDF299,304–305,312,325–328. \textbf{Hȳdan} is printed in Cook, \emph{Exercises in Old English}, V.18.9, p.9 /\allowbreak{} PDF17; the original exercise is unchanged. “Thought; supposed” modernizes Cook's gloss “weened” for \textbf{wēnde}. Dictionary division hyphens in \textbf{gewrit} and \textbf{forlǣtan} are omitted, as elsewhere in the course; their separate practice divisions identify the prefixes.}
% mdparagraph:P100:end


% book-source-body:end
\par\label{end-pronunciation-learner}
\clearpage\endgroup
